% Aufgabenzaehler
\newcounter{exercise}
\newcounter{subexercise}
\newcounter{subexercisenp}
\newcounter{totalexampoints}
\newcounter{exampoints}
\newcounter{number}

\newcommand{\addtoexampoints}[1]{\addtocounter{exampoints}{#1}\addtocounter{exampoints}{-1}\refstepcounter{exampoints}}
\newcommand{\addtototalexampoints}[1]{\addtocounter{totalexampoints}{#1}\addtocounter{totalexampoints}{-1}\refstepcounter{totalexampoints}}
\newcommand{\showtotalpoints}{\IfStrEq{\HidePoints}{ON} {}{von \ref{totalep}}}

\newcommand{\answer}[2]{\IfStrEq{\ANSWER}{ON}{\color{blue}{#1}\color{black}}{#2}}
% Aufgaben werden Automatisch formatiert und nummeriert.
% Aufgabe
\newcommand{\exercise}[5]{\refstepcounter{exercise}\setcounter{exampoints}{0}\setcounter{subexercise}{0}\setcounter{subexercisenp}{0}\vspace{2.5mm}\textbf{\large Aufgabe\, \theexercise~}\IfStrEq{\HidePoints}{OFF}{\ifx\noexpand#2\noexpand\empty  \else \addtoexampoints{#2} \label{\theexercise} \addtototalexampoints{#2}\label{totalep}\fi (\ref{\theexercise}\,Pkt.)}{}\par{#1}\par\vspace{2.5mm}{#3}\answer{#5}{#4}}
% Unteraufgabe
\newcommand{\subexercise}[5]{\refstepcounter{subexercise}\begin{tabu}to \textwidth{X[2,l]X[26,l]X[5,l]}\theexercise.\thesubexercise)&{#1} & \IfStrEq{\HidePoints}{OFF}{~(#2\,Pkt.)}{}\end{tabu}\\{#3}\answer{#5}{#4} \addtoexampoints{#2}\label{\theexercise} \addtototalexampoints{#2}\label{totalep}}
\newcommand{\subexercisenp}[5]{\refstepcounter{subexercisenp}\begin{tabu}to \textwidth{X[1.5,l]X[30,l]}\theexercise.\thesubexercisenp)&{#1}\end{tabu}\\{#3}\answer{#5}{#4}}

% Linien
\newcommand{\letterbox}[2][7mm]{%
  \begin{tikzpicture}[gray]
  \path[use as bounding box](0,0)rectangle(\linewidth,-#2*#1-0.5\pgflinewidth);
   \foreach \n in {1,...,#2}\draw(0 ,-#1*\n )--(\linewidth,-#1*\n );
  \end{tikzpicture}}

% Karierte Flaeche
\newcommand{\plaid}[2][0.5cm]{%
   \begin{tikzpicture}[gray,step=#1]
     \pgfmathtruncatemacro\quantity{(\linewidth-\pgflinewidth)/#1} % maximale Anzahl Kästchen pro Zeile
     \draw (0,0) rectangle (\quantity*#1,#2*#1) (0,0) grid (\quantity*#1,#2*#1);
   \end{tikzpicture}
}

\newcommand{\plaidNotice}{\clearscrheadfoot\ihead{Name:\hspace{105mm}}\ohead{\IfFileExists{latex_exam_lib/bbz-logo.png}{\includegraphics[height=12.5mm]{latex_exam_lib/bbz-logo.png}}{}}\section*{Notizen:}\plaid{39}\clearpage\section*{Notizen:}\plaid{39}}
\newcommand{\lineNotice}{\clearscrheadfoot\ihead{Name:\hspace{105mm}}\ohead{\IfFileExists{latex_exam_lib/bbz-logo.png}{\includegraphics[height=12.5mm]{latex_exam_lib/bbz-logo.png}}{}}\section*{Notizen:}\letterbox{26}\clearpage\section*{Notizen:}\letterbox{26}}
\newcommand{\cloze}[2][10]{\answer{#2\ }{\forloop{number}{0}{\value{number} < #1}{~\_}}}
\newcommand{\oneline}[1]{\begin{tikzpicture}\draw[] node[above]{#1} (-2,0) -- (0.87\textwidth,0);\end{tikzpicture}}
\newcommand{\varline}[4][above]{\begin{tikzpicture}\draw[] node[#1]{#2} (#3,0) -- (#4\textwidth,0);\end{tikzpicture}}
\newcommand{\points}[1]{\textcolor{red}{(#1\,P)}}
\newcommand{\newspaper}[2]{\begin{minipage}{0.495\textwidth}#1\end{minipage}\hspace{0.01\textwidth}\begin{minipage}{0.495\textwidth}#2\end{minipage}}
\newcommand{\leftBlank}{\vfill This page intentionally left blank! \vfill \clearpage}